% !TEX TS-program = pdflatex
% !TeX program = pdflatex
% !TEX encoding = UTF-8
% !TEX spellcheck = fr

\documentclass[xcolor=table]{beamer}


%\usepackage{fullpage}
%\usepackage[left=2.8cm,right=2.2cm,top=2 cm,bottom=2 cm]{geometry}
\setbeamersize{text margin left=10pt,text margin right=10pt}
\usepackage{amsmath,amssymb} 
\usepackage[T1]{fontenc}
\usepackage[utf8]{inputenc}
\usepackage[english,french]{babel}
\usepackage{txfonts}
\usepackage[]{graphicx}
\usepackage{multirow}
\usepackage{hyperref}
\usepackage{colortbl}
\usepackage{listings}
\usepackage{wrapfig}
\usepackage{multicol}
\usepackage[export]{adjustbox} %for images in table, also for frame

\hypersetup{
	colorlinks,
	urlcolor = blue
}

%\renewcommand{\baselinestretch}{1.5}

\def\supit#1{\raisebox{0.8ex}{\small\it #1}\hspace{0.05em}}

\AtBeginSection{%
	\begin{frame}
		\sectionpage
	\end{frame}
}

\newcommand{\rottext}[2]{%
	\rotatebox{90}{%
	\begin{minipage}{#1}%
		\raggedleft#2%
	\end{minipage}%
	}%
}

\usepackage{longtable}
\usepackage{tabu}


\institute{ %
École  nationale Supérieure d'Informatique (ESI, ex. INI), Algérie
}
\author[ \textbf{\footnotesize  \insertframenumber/\inserttotalframenumber} \hspace*{\fill} ESI (2019-2020)] %
{ARIES Abdelkrime}
%\titlegraphic{\includegraphics[height=1cm]{../img/esi-logo.png}%\hspace*{4.75cm}~


\date{Année unniversitaire: 2019/2020} %\today

\usetheme{Warsaw} % Antibes Boadilla Warsaw

\beamertemplatenavigationsymbolsempty

%\setbeamertemplate{headline}{}

\definecolor{lightblue}{HTML}{D0D2FF}
\definecolor{lightyellow}{HTML}{FFFFAA}
\definecolor{darkblue}{HTML}{0000BB}
\definecolor{olivegreen}{HTML}{006600}
\definecolor{violet}{HTML}{6600CC}

\newcommand{\keyword}[1]{\textcolor{red}{\bfseries\itshape #1}}
\newcommand{\expword}[1]{\textcolor{olivegreen}{#1}}
\newcommand{\optword}[1]{\textcolor{violet}{\bfseries #1}}

\makeatletter
\newcommand\mysphere{%
	\parbox[t]{10pt}{\raisebox{0.2pt}{\beamer@usesphere{item projected}{bigsphere}}}}
\makeatother

%\let\oldtabular\tabular
%\let\endoldtabular\endtabular
%\renewenvironment{tabular}{\rowcolors{2}{white}{lightblue}\oldtabular\rowcolor{blue}}{\endoldtabular}


\NoAutoSpacing %french autospacing after ":"

\def\graphpath{}

\newcommand{\changegraphpath}[1]{\def\graphpath{#1}}

\newcommand{\graphintable}[2]{%
	\includegraphics[height=#1,valign=t]{\graphpath #2}%
}

\newcommand{\vgraphpage}[2][.85\textheight]{%
	\begin{center}%
		\includegraphics[height=#1]{\graphpath #2}%
	\end{center}%
}

\newcommand{\hgraphpage}[2][\textwidth]{%
%	\begin{center}%
		\includegraphics[width=#1]{\graphpath #2}%
%	\end{center}%
}

\newcommand{\insertbibliography}[2]{
	\appendix
	\section*{Bibliographie}
	\nocite{#2}
	
	\begin{multicols*}{2}[\frametitle{\insertsection} \usebeamertemplate{frametitle}]%\usebeamertemplate*{frametitle}\frametitle{Références}
		\tiny
		\bibliography{#1}
		\bibliographystyle{acm}
	\end{multicols*}
}



\title[BWEB: 03- Rédaction (Microsoft Word)] %
{Bureautique et Web \\Chapitre 03: Rédaction d'un document numérique \\ \slshape\small  Microsoft Word}  

\changegraphpath{..//img/Bweb03-redaction/word/}
%\setlength{\leftmargini}{2pt}

\settowidth{\leftmargini}{\usebeamertemplate{itemize item}}
\addtolength{\leftmargini}{\labelsep}

\begin{document}
	
	\newcolumntype{L}[2]{%
		>{\vbox to #2\bgroup\vfill\flushleft}%
		p{#1}%
		<{\egroup}} 

\begin{frame}[plain]
	\maketitle
\end{frame}

\begin{frame}
\frametitle{Rédaction d'un document numérique}
\framesubtitle{Microsoft Word}

\begin{itemize}
	\item WYSIWYG: What You See Is What You Get
	\item fait parti de la suite bureautique: Microsoft office
	\item c'est le programme de traitement de texte le plus populaire au monde
	\item facile à configurer et à utiliser 
	\item 
	
\end{itemize}

\end{frame}

\begin{frame}
\frametitle{Rédaction d'un document numérique}
\framesubtitle{Microsoft Word}

%\vgraphpage{google-docs-preview.png}

\end{frame}

\begin{frame}
\frametitle{Rédaction d'un document numérique}
\framesubtitle{Microsoft : Plan}

\begin{multicols}{2}
%	\small
	\tableofcontents
\end{multicols}
\end{frame}

%===================================================================================
\section{Préparer un document}
%===================================================================================

%\begin{frame}
%\frametitle{Préparer un document}
%
%\end{frame}

\subsection{Pages}

\begin{frame}[t]
\frametitle{Préparer un document: Pages}
\framesubtitle{Taille}


\begin{minipage}{0.44\textwidth}
	\begin{itemize}
		\item Onglet: \optword{Mise en page}
		\item Groupe: \optword{Mise en page}
		\item Plusieurs tailles: A4, etc.
		\item On peut définir des tailles personnalisées
	\end{itemize}
\end{minipage}
\begin{minipage}{0.55\textwidth}
	\hgraphpage{taille.png}
\end{minipage}

\end{frame}


\begin{frame}[t]
\frametitle{Préparer un document: Pages}
\framesubtitle{Orientation}

\begin{minipage}{0.44\textwidth}
	\begin{itemize}
		\item Onglet: \optword{Mise en page}
		\item Groupe: \optword{Mise en page}
		\item Deux orientations: \optword{Portrait} ou \optword{Paysage}.
		\item Pour changer l'orientation après une page, insérer un saut de section (page suivante)
	\end{itemize}
\end{minipage}
\begin{minipage}{0.55\textwidth}
	\hgraphpage{orientation.png}
\end{minipage}

\end{frame}

\begin{frame}[t]
\frametitle{Préparer un document: Pages}
\framesubtitle{Marges}

\begin{minipage}{0.44\textwidth}
	\begin{itemize}
		\item Onglet: \optword{Mise en page}
		\item Groupe: \optword{Mise en page}
		\item Des marges prédéfinies
		\item On peut ajouter des marges personnalisées
	\end{itemize}
\end{minipage}
\begin{minipage}{0.55\textwidth}
	\hgraphpage{marges.png}
\end{minipage}

\end{frame}

\begin{frame}[t]
\frametitle{Préparer un document: Pages}
\framesubtitle{Colonnes}

\begin{minipage}{0.44\textwidth}
	\begin{itemize}
		\item Onglet: \optword{Mise en page}
		\item Groupe: \optword{Mise en page}
		\item Sélectionner le texte qu'on veuille mettre en colonnes
		\item Choisir le nombre de colonnes 
		\item Dans les articles scientifiques, la double-colonnes est très utilisée
	\end{itemize}
\end{minipage}
\begin{minipage}{0.55\textwidth}
	\hgraphpage{colonnes.png}
\end{minipage}

\end{frame}

\begin{frame}[t]
\frametitle{Préparer un document: Pages}
\framesubtitle{Saute de page et section}

\begin{minipage}{0.44\textwidth}
	\begin{itemize}
		\item Onglet: \optword{Mise en page}
		\item Groupe: \optword{Mise en page}
		\item On insère un saut de page par exemple après la table de matière, liste des tableaux, etc.
		\item On insère un saut de section pour séparer les chapitres: avoir des différentes entêtes, différentes numérotations, etc.
	\end{itemize}
\end{minipage}
\begin{minipage}{0.55\textwidth}
	\hgraphpage{sautes.png}
\end{minipage}

\end{frame}

\begin{frame}[t]
\frametitle{Préparer un document: Pages}
\framesubtitle{Page de garde}

\begin{minipage}{0.49\textwidth}
	\begin{itemize}
		\item Onglet: \optword{Insertion}
		\item Groupe: \optword{Pages}
		\item On peut choisir une page de garde
		\item ou, naviguer d'autre sur le site \url{office.com}
	\end{itemize}
\end{minipage}
\begin{minipage}{0.50\textwidth}
	\hgraphpage{garde.png}
\end{minipage}

\end{frame}

\begin{frame}[t]
\frametitle{Préparer un document: Pages}
\framesubtitle{Bordures}

\begin{minipage}{0.49\textwidth}
	\begin{itemize}
		\item Onglet: \optword{Création}
		\item Groupe: \optword{Arrière-plan de la page}
		\item On peut ajouter des bordures de page:
		\begin{itemize}
			\item Encadrement
			\item Ombre
			\item 3D
		\end{itemize}
	\end{itemize}
\end{minipage}
\begin{minipage}{0.50\textwidth}
	\hgraphpage{bordures.png}
\end{minipage}

\end{frame}

\begin{frame}[t]
\frametitle{Préparer un document: Pages}
\framesubtitle{Filigrane}

\begin{minipage}{0.49\textwidth}
	\begin{itemize}
		\item Onglet: \optword{Création}
		\item Groupe: \optword{Arrière-plan de la page}
		\item Le filigrane (Watermark) est un texte transparent qui apparait sur les pages
		\item On peut: 
		\begin{itemize}
			\item choisir un filigrane prédéfinie
			\item naviguer d'autres sur le site \url{office.com}
			\item supprimer le filigrane actuel 
			\item créer un filigrane personnalisé
		\end{itemize}
	\end{itemize}
\end{minipage}
\begin{minipage}{0.50\textwidth}
	\hgraphpage{filigrane1.png}
\end{minipage}

\end{frame}

\begin{frame}[t]
\frametitle{Préparer un document: Pages}
\framesubtitle{Filigrane}

\begin{minipage}{0.49\textwidth}
	\begin{itemize}
		\item Onglet: \optword{Création}
		\item Groupe: \optword{Arrière-plan de la page}
		\item On peut choisir une image ou un texte 
	\end{itemize}
\end{minipage}
\begin{minipage}{0.50\textwidth}
	\hgraphpage{filigrane2.png}
\end{minipage}

\end{frame}

\begin{frame}[t]
\frametitle{Préparer un document: Pages}
\framesubtitle{Arrière-plan}

\begin{minipage}{0.44\textwidth}
	\begin{itemize}
		\item Onglet: \optword{Création}
		\item Groupe: \optword{Arrière-plan de la page}
		\item Pour choisir la couleur des pages
		\item on peut choisir d'autres couleurs 
		\item ou, ajouter un effet de remplissage
	\end{itemize}
\end{minipage}
\begin{minipage}{0.55\textwidth}
	\hgraphpage{couleur.png}
\end{minipage}

\end{frame}

\begin{frame}[t]
\frametitle{Préparer un document: Pages}
\framesubtitle{Arrière-plan}

\begin{minipage}{0.44\textwidth}
	\begin{itemize}
		\item Onglet: \optword{Création}
		\item Groupe: \optword{Arrière-plan de la page}
		\item Un effet peut être: 
		\begin{itemize}
			\item Dégradé : 
			\item Texture : 
			\item Motif : 
			\item Image : 
		\end{itemize}
	\end{itemize}
\end{minipage}
\begin{minipage}{0.55\textwidth}
	\hgraphpage{remplissage.png}
\end{minipage}

\end{frame}



\begin{frame}
\frametitle{Préparer un document: Pages}
\framesubtitle{En-tête et pied de page}

\begin{minipage}{0.78\textwidth}
	\begin{itemize}
		\item Onglet: \optword{Insertion} 
		\item Groupe: \optword{En-tête et pied de page}
	\end{itemize}
\end{minipage}
\begin{minipage}{0.2\textwidth}
	\hgraphpage{entete-pied.png}
\end{minipage}

\vspace{6pt}
Si on double-clique sur l'entête, un onglet va s'apparaitre
\begin{itemize}
	\item \optword{Première page différente}: en générale pour la page de garde 
	\item \optword{Pages paires et impaires différentes}: par exemple, pour afficher le nom du chapitre dans une page, et le nom de l'auteur dans la suivante
	\item \optword{Lier au précédent}: après un saut de section, on peut changer le style des entêtes et pieds de page en découchant cette option
\end{itemize}

\hgraphpage{entete-pied-barre.png}

\end{frame}


\subsection{Paragraphes}

\begin{frame}
\frametitle{Préparer un document: Paragraphes}
\framesubtitle{Paragraphes}

\begin{minipage}{0.38\textwidth}
	\begin{itemize}
		\item Onglet: \optword{Accueil} et \optword{Mise en page}
		\item Groupe: \optword{Paragraphe}
		\item Fonctionnalités:
		\begin{itemize}
			\item Alignement 
			\item Espacement vertical
			\item Interligne
			\item Retrait
			\item Retrait de 1ère ligne
			\item Listes
			\item Affichage des symboles non imprimables
		\end{itemize}
	\end{itemize}
\end{minipage}
\begin{minipage}{0.60\textwidth}
	\def\arraystretch{.5}
	\begin{tabular}{@{}ll}
		\hgraphpage[.48\textwidth]{accueil-paragraphe.png} &
		\hgraphpage[.48\textwidth]{misepage-paragraphe.png} \\
		\tiny Accueil & \tiny Mise en page \\
	\end{tabular}
%	\begin{center}
%		\vspace{-8pt}
		\hgraphpage[.62\textwidth]{paragraphe.png}
%	\end{center}
\end{minipage}

\end{frame}

\subsection{Mise en forme et styles}

\begin{frame}
\frametitle{Préparer un document: Mise en forme et styles}
\framesubtitle{Mise en forme}

\begin{minipage}{0.48\textwidth}
	\begin{itemize}
		\item Police
		\item Couleur 
		\item Casse 
	\end{itemize}
\end{minipage}
\begin{minipage}{0.50\textwidth}
	\hgraphpage{miseforme.png}
	
	\hgraphpage[0.8\textwidth]{police.png}
	%	\end{center}
\end{minipage}

\end{frame}

\begin{frame}
\frametitle{Préparer un document: Mise en forme et styles}
\framesubtitle{Casse et effets de texte}

\begin{minipage}{0.58\textwidth}
	\begin{itemize}
		\item Police
		\item Couleur 
		\item Casse 
	\end{itemize}
\end{minipage}
\begin{minipage}{0.4\textwidth}
	\hgraphpage[.5\textwidth]{casse.png}
	
	\hgraphpage{effet-texte.png}
	%	\end{center}
\end{minipage}

\end{frame}


\begin{frame}
\frametitle{Préparer un document: Mise en forme et styles}
\framesubtitle{Styles}

\begin{minipage}{0.58\textwidth}
	\begin{itemize}
		\item Police
		\item Couleur 
		\item Casse 
	\end{itemize}
\end{minipage}
\begin{minipage}{0.4\textwidth}
	\hgraphpage{style.png}
	
	\vspace{12pt}
	
	\hgraphpage[.8\textwidth]{style-new.png}
\end{minipage}

\end{frame}

\subsection{Modification}%Recherche et remplacement

\begin{frame}[t]
\frametitle{Préparer un document: Modification}
\framesubtitle{Recherche}

\begin{minipage}{0.60\textwidth}
	\begin{itemize}
		\item 
	\end{itemize}
\end{minipage}
\begin{minipage}{0.38\textwidth}
	\begin{flushright}
		\hgraphpage[1.5cm]{edition.png}
	\end{flushright}

	\vspace{-12pt}
	\hgraphpage{rechercher.png}
\end{minipage}

\end{frame}

\begin{frame}[t]
\frametitle{Préparer un document: Modification}
\framesubtitle{Remplacement}

\begin{minipage}{0.38\textwidth}
	\begin{itemize}
		\item 
	\end{itemize}
\end{minipage}
\begin{minipage}{0.6\textwidth}	
	\begin{flushright}
		\hgraphpage[1.5cm]{edition.png}
	\end{flushright}

	\vspace{-12pt}
	\hgraphpage{remplacer.png}
\end{minipage}

\end{frame}

%===================================================================================
\section{Enrichir un document}
%===================================================================================

%\begin{frame}
%\frametitle{Enrichir un document}
%
%\end{frame}

\subsection{Tableaux}


\begin{frame}
\frametitle{Enrichir un document: Tableaux}
\framesubtitle{Création}

\begin{minipage}{0.70\textwidth}
	\begin{itemize}
		\item 
	\end{itemize}
\end{minipage}
\begin{minipage}{0.28\textwidth}	
	\hgraphpage{tableau.png}
\end{minipage}

\end{frame}

\begin{frame}[t]
\frametitle{Enrichir un document: Tableaux}
\framesubtitle{Options: Disposition}

\hgraphpage{tableau-disposition.png}

\end{frame}

\begin{frame}[t]
\frametitle{Enrichir un document: Tableaux}
\framesubtitle{Options: Disposition}

\hgraphpage{tableau-disposition.png}

\end{frame}


\subsection{Illustrations}

\begin{frame}[t]
\frametitle{Enrichir un document: Illustrations}
\framesubtitle{Insérer une illustration}

\begin{minipage}{0.43\textwidth}
	\begin{itemize}
		\item 
	\end{itemize}
\end{minipage}
\begin{minipage}{0.28\textwidth}
	\hgraphpage{illustrations.png}
	
	\hgraphpage{illustrations-graphique.png}
	
	\hgraphpage{illustrations-smartart.png}
\end{minipage}
\begin{minipage}{0.27\textwidth}	
	\hgraphpage{illustrations-formes.png}
\end{minipage}


\end{frame}


\begin{frame}[t]
\frametitle{Enrichir un document: Illustrations}
\framesubtitle{Outils image}

\hgraphpage{illustration-image-outils.png}

\begin{minipage}{0.56\textwidth}
	\begin{itemize}
		\item 
	\end{itemize}
\end{minipage}
\begin{minipage}{0.22\textwidth}
	\hgraphpage{illustrations-habillage.png}
\end{minipage}
\begin{minipage}{0.20\textwidth}	
	\hgraphpage{illustrations-position.png}
\end{minipage}

\end{frame}

\subsection{Liens et renvoi}

\begin{frame}
\frametitle{Enrichir un document: Liens et renvoi}
\framesubtitle{Insérer un lien hypertexte}

\end{frame}

\begin{frame}
\frametitle{Enrichir un document: Liens et renvoi}
\framesubtitle{Renvoi vers un emplacement dans le document}

\end{frame}

\subsection{Formules et symboles}

\begin{frame}
\frametitle{Enrichir un document: Formules et symboles}
\framesubtitle{Insérer une formule}

\hgraphpage{equation-barre.png}

\begin{minipage}{0.76\textwidth}
	\begin{itemize}
		\item 
	\end{itemize}
\end{minipage}
\begin{minipage}{0.23\textwidth}
	\hgraphpage{equation.png}
		
	\hgraphpage{equation-manuelle.png}
\end{minipage}

\end{frame}

\begin{frame}
\frametitle{Enrichir un document: Formules et symboles}
\framesubtitle{Insérer des symboles}

\begin{minipage}{0.50\textwidth}
	\begin{itemize}
		\item 
	\end{itemize}
\end{minipage}
\begin{minipage}{0.49\textwidth}
	\hgraphpage[.3\textwidth]{symboles.png}
	
	\hgraphpage{symboles-table.png}
\end{minipage}

\end{frame}

%===================================================================================
\section{Références}
%===================================================================================

%\begin{frame}
%\frametitle{Références}
%
%\end{frame}

\subsection{Notes de bas de page}

\begin{frame}[t]
\frametitle{Références: Notes de bas de page}
\framesubtitle{Insérer une note de bas de page}

\begin{minipage}{0.69\textwidth}
	\begin{itemize}
		\item 
	\end{itemize}
\end{minipage}
\begin{minipage}{0.30\textwidth}
	\hgraphpage{bas-page.png}
\end{minipage}

\begin{itemize}
	\item 
\end{itemize}

\end{frame}

\subsection{Table de matières}

\begin{frame}
\frametitle{Références: Table de matières}
\framesubtitle{Insérer une table de matières}

\begin{minipage}{0.59\textwidth}
	\begin{itemize}
		\item Onglet: \optword{Références}
		\item Groupe: \optword{Légendes}
		\item Mettre le curseur dans la position où on veut insérer la table
		\item Cliquer sur \optword{Insérer une table des illustrations}
		\item Choisir le type de la légende: Figure ou Tableau
		\item Choisir le style
	\end{itemize}
\end{minipage}
\begin{minipage}{0.40\textwidth}
	\hgraphpage{table-matiere.png}
\end{minipage}

\end{frame}

\begin{frame}
\frametitle{Références: Table de matières}
\framesubtitle{Insérer une table de matières personnalisée}

\begin{minipage}{0.59\textwidth}
	\begin{itemize}
		\item 
	\end{itemize}
\end{minipage}
\begin{minipage}{0.40\textwidth}
	\hgraphpage{table-matiere-perso.png}
	
	\hgraphpage[.6\textwidth]{table-matiere-options.png}
\end{minipage}

\end{frame}

\subsection{Table des illustrations}

\begin{frame}[t]
\frametitle{Références: Table des illustrations}
\framesubtitle{Insérer des légendes}

\begin{minipage}{0.61\textwidth}
	\begin{itemize}
		\item Onglet: \optword{Références}
		\item Groupe: \optword{Légendes}
		\item Sélectionner l'illustration (tableau ou image)
		\item Cliquer sur \optword{Insérer une légende}
		\item Choisir le type de la légende: Figure ou Tableau
		\item Ajouter un titre à la légende
		\item Pour changer le style de numérotation, appuyer sur \optword{Numérotation}
	\end{itemize}
\end{minipage}
\begin{minipage}{0.38\textwidth}
	\hgraphpage[.6\textwidth]{legende-barre.png}
	
	\hgraphpage{legende.png}
	
	\hgraphpage[.8\textwidth]{legende-num.png}
\end{minipage}

\end{frame}

\begin{frame}[t]
\frametitle{Références: Table des illustrations}
\framesubtitle{Insérer une table des illustrations}

\begin{minipage}{0.61\textwidth}
	\begin{itemize}
		\item Onglet: \optword{Références}
		\item Groupe: \optword{Légendes}
		\item Mettre le curseur dans la position où on veut insérer la table
		\item Cliquer sur \optword{Insérer une table des illustrations}
		\item Choisir le type de la légende: Figure ou Tableau
		\item Choisir le style
	\end{itemize}
\end{minipage}
\begin{minipage}{0.38\textwidth}
	\hgraphpage[.6\textwidth]{legende-barre.png}
	
	\hgraphpage{table-illustrations.png}
	
\end{minipage}

\end{frame}

\subsection{Index}

\begin{frame}[t]
\frametitle{Références: Index}
\framesubtitle{Insérer des entrées}

\begin{minipage}{0.61\textwidth}
\begin{itemize}
	\item 
\end{itemize}
\end{minipage}
\begin{minipage}{0.38\textwidth}
\hgraphpage[.6\textwidth]{index-barre.png}

\hgraphpage[.8\textwidth]{entree.png}

\end{minipage}

\end{frame}

\begin{frame}[t]
\frametitle{Références: Table des illustrations}
\framesubtitle{Insérer un index}

\begin{minipage}{0.61\textwidth}
\begin{itemize}
\item 
\end{itemize}
\end{minipage}
\begin{minipage}{0.38\textwidth}
\hgraphpage[.6\textwidth]{index-barre.png}

\hgraphpage{index.png}

\end{minipage}

\end{frame}

%===================================================================================
\section{Révision et partage}
%===================================================================================

%\begin{frame}
%\frametitle{Révision et partage}
%
%\end{frame}

\subsection{Vérification}

\begin{frame}
\frametitle{Révision et partage}
\framesubtitle{Vérification}

\end{frame}

\subsection{Commentaires et suivi}

\begin{frame}
\frametitle{Révision et partage}
\framesubtitle{Commentaires et suivi}

\end{frame}

\subsection{Comparaison}

\begin{frame}
\frametitle{Révision et partage}
\framesubtitle{Comparaison}

\end{frame}

\subsection{Partage}

\begin{frame}
\frametitle{Révision et partage}
\framesubtitle{Partage}


\end{frame}

\begin{frame}
\frametitle{Révision et partage: Partage}
\framesubtitle{Protection}

\end{frame}

\begin{frame}
\frametitle{Révision et partage: Partage}
\framesubtitle{Impression}

\end{frame}

\begin{frame}
\frametitle{Révision et partage: Partage}
\framesubtitle{Sauvegarde}

\end{frame}

\begin{frame}
\frametitle{Révision et partage: Partage}
\framesubtitle{Sauvegarde sur cloud}
\end{frame}

\begin{frame}
\frametitle{Rédaction d'un document numérique: Microsoft Word}
\framesubtitle{Un peu d'humeur}

\hgraphpage{word-humor.jpg}

\end{frame}

\insertbibliography{Bweb03}{*}

%\appendix
%\section*{Bibliographie}
%\nocite{*}
%
%\begin{multicols*}{2}[\frametitle{\insertsection} \usebeamertemplate{frametitle}]%\usebeamertemplate*{frametitle}\frametitle{Références}
%	\tiny
%	\bibliography{Bweb03}
%	\bibliographystyle{acm}
%\end{multicols*}


\end{document}

